% Figure: classification of critical points by Hessian eigenvalues.
% Four panels: positive-definite (bowl), negative-definite (dome),
% indefinite (saddle), semi-definite (degenerate trough).
\begin{figure}[htbp]
\centering
\begin{tikzpicture}[scale=0.8,
  contour/.style={engblue, thick},
  panel/.style={draw=engblack, thick},
  caption/.style={font=\small, align=center},
  star/.style={engred},
]

% ---- Panel 1: positive definite (local min): f = x^2 + y^2 ----
\begin{scope}[shift={(0, 0)}]
  \draw[panel] (-1.6, -1.6) rectangle (1.6, 1.6);
  \foreach \r in {0.4, 0.8, 1.2} {
    \draw[contour] (0, 0) circle (\r);
  }
  \filldraw[star] (0,0) circle (0.06);
  \node[caption] at (0, -2.05) {min\\$\lambda_{1}, \lambda_{2} > 0$};
\end{scope}

% ---- Panel 2: negative definite (local max): f = -x^2 - y^2 ----
\begin{scope}[shift={(4, 0)}]
  \draw[panel] (-1.6, -1.6) rectangle (1.6, 1.6);
  \foreach \r in {0.4, 0.8, 1.2} {
    \draw[contour] (0, 0) circle (\r);
  }
  \filldraw[star] (0,0) circle (0.06);
  \node[caption] at (0, -2.05) {max\\$\lambda_{1}, \lambda_{2} < 0$};
\end{scope}

% ---- Panel 3: indefinite (saddle): f = x^2 - y^2 ----
\begin{scope}[shift={(8, 0)}]
  \draw[panel] (-1.6, -1.6) rectangle (1.6, 1.6);
  % Hyperbolae x^2 - y^2 = c for c = +/- 0.25, +/- 0.81, +/- 1.44
  % Restrict domain to real-valued branches: |x| >= sqrt(c) for the y(x) branches
  \foreach \c [evaluate=\c as \r using {sqrt(\c)}] in {0.25, 0.81, 1.44} {
    \draw[contour, domain=\r:1.55, samples=40]
      plot ({\x}, {sqrt(\x*\x - \c)});
    \draw[contour, domain=\r:1.55, samples=40]
      plot ({\x}, {-sqrt(\x*\x - \c)});
    \draw[contour, domain=-1.55:-\r, samples=40]
      plot ({\x}, {sqrt(\x*\x - \c)});
    \draw[contour, domain=-1.55:-\r, samples=40]
      plot ({\x}, {-sqrt(\x*\x - \c)});
    \draw[contour, domain=-1.55:1.55, samples=80]
      plot ({sqrt(\x*\x + \c)}, {\x});
    \draw[contour, domain=-1.55:1.55, samples=80]
      plot ({-sqrt(\x*\x + \c)}, {\x});
  }
  % Asymptotes
  \draw[engblack!50] (-1.5, -1.5) -- (1.5, 1.5);
  \draw[engblack!50] (-1.5, 1.5) -- (1.5, -1.5);
  \filldraw[star] (0,0) circle (0.06);
  \node[caption] at (0, -2.05) {saddle\\$\lambda_{1} > 0,\ \lambda_{2} < 0$};
\end{scope}

% ---- Panel 4: semi-definite (degenerate trough): f = x^2 ----
\begin{scope}[shift={(12, 0)}]
  \draw[panel] (-1.6, -1.6) rectangle (1.6, 1.6);
  % Vertical lines x = +/- sqrt(c) for c = 0.25, 1, 1.96 (but 1.96 > 1.6 so skip)
  \foreach \c in {0.25, 0.81, 1.44} {
    \draw[contour] ({sqrt(\c)}, -1.55) -- ({sqrt(\c)}, 1.55);
    \draw[contour] ({-sqrt(\c)}, -1.55) -- ({-sqrt(\c)}, 1.55);
  }
  \draw[engblack!60, thick, dashed] (0, -1.55) -- (0, 1.55);
  \filldraw[star] (0,0) circle (0.06);
  \node[caption] at (0, -2.05) {degenerate\\$\lambda_{1} > 0,\ \lambda_{2} = 0$};
\end{scope}

\end{tikzpicture}
\caption[Critical points classified by the Hessian
eigenvalues.]{Four canonical critical-point types in two variables,
shown as contour plots near the origin. From left: a local minimum
(both Hessian eigenvalues positive, contours nested closed curves
around the critical point); a local maximum (both eigenvalues
negative, geometrically indistinguishable from the minimum without
the sign of the function values); a saddle (one eigenvalue of each
sign, contours hyperbolae with the critical point at the asymptote
intersection); a degenerate critical point (one zero eigenvalue,
contours degenerate to parallel lines, the second-derivative test
inconclusive). The red dot marks the critical point in each panel.}
\label{fig:vol02:ch11:critical-classification}
\end{figure}
